This paper studies how benchmark scores of automatically tuned LLM systems should be interpreted and reported. We expect a positive effect on scientific practice (less over-claiming, fairer comparisons between systems tuned with different budgets). We see no direct negative societal consequences; one risk is that the more conservative intervals we advocate may be read as a general reason to distrust any tuned system, whereas our results are specific to the regimes studied. All experiments use public data and small open-weight models on CPU.
